% GEDES defense reviewer — 50 questions with worked answers.
% Build:  latexmk -pdf DEFENSE_REVIEWER.tex
% Figures are pulled from the real thesis/guidebook sources; nothing here is
% redrawn for effect. Every number is traceable to code/results/*.json.
\documentclass[11pt,a4paper]{article}

\usepackage[T1]{fontenc}
\usepackage[utf8]{inputenc}
\usepackage{lmodern}
\usepackage[margin=2.3cm,top=2.0cm,bottom=2.2cm]{geometry}
\usepackage{amsmath,amssymb}
\usepackage{graphicx}
\usepackage{booktabs}
\usepackage{xcolor}
\usepackage[most]{tcolorbox}
\usepackage{enumitem}
\usepackage{titlesec}
\usepackage[hidelinks]{hyperref}

\graphicspath{{img/}{../../../figures/}{../../jgr/guidebook/figures-tikz/}}

\definecolor{char}{HTML}{2A2520}
\definecolor{coral}{HTML}{B85B3A}
\definecolor{teal}{HTML}{2A6478}
\definecolor{forest}{HTML}{3D6E4A}
\definecolor{dim}{HTML}{6E6862}
\definecolor{grid}{HTML}{E8E5E0}
\definecolor{tint}{HTML}{F7F5F2}

\color{char}
\setlength{\parindent}{0pt}
\setlength{\parskip}{5pt}

\titleformat{\section}{\Large\bfseries\color{teal}}{\thesection}{0.6em}{}
\titlespacing{\section}{0pt}{20pt}{8pt}

\newcounter{qq}
\newcommand{\Q}[1]{%
  \refstepcounter{qq}%
  \vspace{12pt}%
  \noindent{\color{coral}\bfseries Q\theqq.}\enspace{\bfseries #1}\par
  \vspace{3pt}}

\newtcolorbox{say}{colback=forest!7,colframe=forest,boxrule=0.9pt,
  arc=2pt,left=7pt,right=7pt,top=5pt,bottom=5pt,
  fonttitle=\bfseries\small\color{forest},title={SAY THIS}}

\newtcolorbox{trap}{colback=coral!7,colframe=coral,boxrule=0.9pt,
  arc=2pt,left=7pt,right=7pt,top=5pt,bottom=5pt,
  fonttitle=\bfseries\small\color{coral},title={DO NOT SAY}}

\newcommand{\pointat}[1]{\par\vspace{2pt}%
  {\footnotesize\color{dim}\textit{Point at: #1}}\par}
\newcommand{\num}[1]{\textbf{#1}}

\begin{document}

\thispagestyle{empty}
\begin{center}
{\Huge\bfseries Defense reviewer}\\[6pt]
{\large 51 questions, with the answers and the numbers behind them}\\[16pt]
{\color{dim}Ramon III P. Gregorio\quad$\cdot$\quad 24M58378\quad$\cdot$\quad
Institute of Science Tokyo, Kasai Laboratory}\\[4pt]
{\color{dim}\small Difference of Lunar Regolith Thermal Conductivity $K_d$ at the
Apollo 15 and 17 Heat-Flow Boreholes}
\end{center}

\vspace{10pt}
\hrule
\vspace{10pt}

\textbf{How to use this.} Each question has a \textcolor{forest}{SAY THIS} box:
that is the 20--40 second spoken answer, and it is the only part you have to
have ready word-for-word. Everything under it is the backing you need if the
examiner pushes. \textcolor{coral}{DO NOT SAY} boxes mark answers that sound
reasonable and will lose you the room.

\textbf{Every number here is traceable.} Sources are
\texttt{code/results/*.json} and \texttt{code/src/lunar/config.py}. If a number
in this document disagrees with the code, the code is right.

\vspace{6pt}
\begin{tcolorbox}[colback=tint,colframe=grid,boxrule=0.8pt,arc=2pt]
{\bfseries The twelve numbers you must know cold.}\par\vspace{4pt}
\small
\begin{tabular}{@{}ll@{\hskip 24pt}ll@{}}
$K_d^{*}$ Apollo 15 & \num{4.60} mW m$^{-1}$ K$^{-1}$ &
  95\,\% CI A15 & \num{[4.18, 6.96]} \\
$K_d^{*}$ Apollo 17 & \num{7.08} &
  95\,\% CI A17 & \num{[6.16, 8.07]} \\
Global (Hayne 2017) & \num{3.4} &
  Contrast & \num{2.31}, CI \num{[-0.12, 3.56]} \\
Deep sensors & \num{7} (A15), \num{16} (A17) &
  Bootstrap tail & \num{0.031} = 46/1500 \\
Basal flux $Q_b$ & \num{21} / \num{16} mW m$^{-2}$ &
  RMSE A17, global$\to$fit & \num{0.89} $\to$ \num{0.40} K \\
Anchor depth & \num{0.55} m &
  Speed-up & $\approx$\num{2500}$\times$, 27 h $\to$ $<$1 min \\
\end{tabular}
\end{tcolorbox}

\vspace{4pt}
{\footnotesize\color{dim}Timing: 12 min talk + 5 min Q\&A on the thesis;
6 min + 7 min on the doctoral plan. \textbf{The doctoral Q\&A is longer than the
thesis Q\&A} --- Part F matters more than its length here suggests.}

\clearpage

% =====================================================================
\section{A \; The problem and the data}

\Q{In one sentence: what did you do, and what is new?}
\begin{say}
Every thermal model of the Moon uses a single conductivity for the whole body.
I tested that assumption at the only two places it can be tested --- the Apollo
15 and 17 heat-flow boreholes --- and found that both sites require a higher
value than the global one, and that they differ from each other. To make that
test affordable I built a solver that returns the same answer about 2500 times
faster.
\end{say}

There are two contributions and you should always name both, in this order:
the \emph{measurement} (per-site $K_d$, and the contrast between the sites) and
the \emph{method} (the flux-anchored equilibrium solver). The method is what
turned a one-day calculation into a one-minute one, which is what made a
1500-draw uncertainty analysis possible at all. Without the method there is no
error bar; without the error bar there is no defensible claim.

\Q{Why does the deep temperature matter? Nobody lives a metre underground.}
\begin{say}
Three reasons. Water ice can only survive where the ground is cold and steady,
and that is decided below the daily heat wave, not at the surface. Instruments
that see beneath the surface --- microwave and terahertz sounders --- need a
subsurface temperature profile as a boundary condition, and no satellite can
measure one. And the deep temperature is set by the conductivity, which is the
property I retrieve.
\end{say}

The surface swings roughly 100--390 K every lunation. By about 15--20 cm that
swing is essentially gone, and what remains is a slowly rising mean profile set
by the geothermal flux divided by the conductivity. That steady value is the
quantity ice stability and radiative-transfer retrievals actually depend on.

\pointat{Slide 2, the animation --- the surface readout never stops moving and
the one-metre readout never starts.}

\Q{Where does the global value 3.4 come from, and why do you say it is unchecked?}
\begin{say}
It comes from Hayne et al.\ 2017, fitted to Diviner orbital infrared data. The
issue is not that it was done carelessly --- it is that orbital infrared only
senses the top few centimetres, the part that is swinging wildly. The value has
never been compared against a measurement from below the diurnal skin, because
until this work nobody had done the comparison at the only two sites where the
data exists.
\end{say}

\begin{trap}
``The Hayne value is wrong.'' It is not wrong; it is a well-constrained fit to
a different quantity, in a different depth range, for a different purpose. Your
claim is narrower and stronger: it has never been tested at depth, and when you
test it at depth, both boreholes want something higher.
\end{trap}

\Q{Why only two sites? Isn't that too small a sample to say anything?}
\begin{say}
It is the entire population, not a sample. Apollo 15 and 17 are the only places
where humans have ever drilled into another body and left thermometers. I am
not claiming these two sites represent the Moon --- I am claiming that the one
number applied to the whole Moon fails at both places where it can be checked,
and that is a meaningful statement about the assumption even with $n=2$.
\end{say}

This is the honest framing and it is stronger than a claim of generality. Two
out of two is not a statistical population but it is a complete audit of the
available evidence. The doctoral plan exists precisely because two points are
not a map.

\Q{What exactly is in the Apollo dataset?}
\begin{say}
Two probes per site. Apollo 15 reaches 1.4 m, Apollo 17 reaches 2.3 m. Each
probe carries platinum-resistance thermometers in bridge pairs, recording from
1971 to 1977. After the 80 cm exclusion I have 23 usable sensors: 7 at Apollo
15 and 16 at Apollo 17. The record I use is the restored one recovered from the
original mission tapes by Nagihara and colleagues in 2018.
\end{say}

Points worth having ready: the probes sit inside a fibreglass--epoxy
\emph{borestem} left in the hole; sensor depths carry roughly $\pm2.5$ cm
emplacement uncertainty, which is exactly the jitter used in the bootstrap; and
the Apollo 17 probes are the deeper pair, which is why that site carries most
of the constraining power.

\begin{trap}
Do not say the probe is aluminium. It is a fibreglass--epoxy borestem. This is
a favourite check question.
\end{trap}

\pointat{Slide 4, both boreholes to scale with the Clementine globe.}

\Q{Why do you throw away every sensor shallower than 80 cm?}
\begin{say}
Because of drilling disturbance, not because of the diurnal wave. Emplacing the
probe disturbed the regolith around the hole, and the borestem itself
short-circuits heat down the hole. The upper sensors are contaminated by the
instrument, so they are excluded. I then checked that the result does not
depend on where exactly I put the cut.
\end{say}

\begin{trap}
``80 cm is where the daily temperature swing becomes negligible.'' The diurnal
wave is gone by roughly 15--20 cm; the thermal skin depth is about 5 cm. An
examiner who knows that will take the slide apart. 80 cm is a
\emph{contamination} criterion.
\end{trap}

Then quote the sensitivity, from \texttt{borestem\_sensitivity.json}:

\vspace{4pt}
{\small
\begin{tabular}{@{}lcccc@{}}
\toprule
cut & 60 cm & 70 cm & \textbf{80 cm} & 90 cm \\
\midrule
Apollo 15 & 4.79 ($n{=}8$) & 4.79 ($n{=}8$) & \textbf{4.60} ($n{=}7$) & 4.83 ($n{=}5$) \\
Apollo 17 & 8.34 ($n{=}18$) & 7.08 ($n{=}16$) & \textbf{7.08} ($n{=}16$) & 7.08 ($n{=}16$) \\
contrast  & 3.56 & 2.29 & \textbf{2.48} & 2.25 \\
\bottomrule
\end{tabular}}
\vspace{4pt}

Apollo 17 is \emph{identical} for every cut at or below 70 cm; only at 60 cm do
two contaminated sensors drag it up to 8.34. Apollo 15 moves by about
$\pm0.12$. The ordering survives every choice. This turns your most
discretionary decision into a robustness result.

\Q{How do you turn six years of drifting record into one temperature per sensor?}
\begin{say}
With an automatic rule, not by eye. For each sensor I scan candidate window
starts across the last 55 to 85 percent of the record, fit a straight line to
the trailing part, and keep the earliest start whose drift is below 0.08 kelvin
per year. That gives the longest clean window, which gives the smallest
statistical error. The equilibrium temperature is the mean over that window,
and any leftover drift is carried into the error budget rather than discarded.
\end{say}

If a sensor has no qualifying window, the rule falls back to the final quarter
of the record and flags it --- it does not silently pass. The rule is
reproducible: the same code runs on every sensor and anyone can re-run it.

Robustness, from \texttt{stability\_threshold\_sensitivity.json}: moving the
drift threshold from 0.04 to 0.16 K yr$^{-1}$ changes $K_d^{*}$ by less than
0.02 mW at either site, and never changes which sensors qualify.

\pointat{Slide 6, and backup: the guidebook \texttt{windowflow} chart.}

% =====================================================================
\section{B \; The physics and the model}

\Q{Write down the equation you are solving and name every term.}
\begin{say}
One-dimensional heat conduction in a vertical column of regolith:
$\rho(z)\,c_p(T)\,\partial T/\partial t = \partial_z[K(T,z)\,\partial_z T]$.
Density varies with depth, heat capacity with temperature, and the conductivity
with both. The only unknown is the deep value of the conductivity.
\end{say}

Everything else is measured or published: $\rho(z)$ and $c_p(T)$ from Hayne
2017 and Hemingway 1973, the forcing from SPICE ephemeris geometry and the
site albedo, and $Q_b$ from Langseth 1976.

\Q{What is $K(T,z)$, and where does each constant come from?}
\begin{say}
It is the Hayne 2017 two-part form: a contact term that rises with depth as the
regolith compacts, multiplied by a radiative term that grows as temperature
cubed because heat also crosses the pore space by radiation.
\end{say}
\[
K(T,z)=\underbrace{\bigl[K_d-(K_d-K_s)\,e^{-z/H}\bigr]}_{\text{contact, compaction}}
\times
\underbrace{\bigl[1+\chi\,(T/T_{\rm ref})^{3}\bigr]}_{\text{radiative}}
\]
with $K_s=7.4\times10^{-4}$ W m$^{-1}$ K$^{-1}$, $H=6$ cm, $\chi=2.7$,
$T_{\rm ref}=350$ K --- all fixed at their published values. $K_d$ is the only
free parameter in the entire retrieval.

\Q{What is $\chi$, and why is it your largest caveat?}
\begin{say}
$\chi$ sets how much heat crosses the pore space by radiation rather than by
grain contact. I hold it at Hayne's 2.7. It is my largest caveat because my
absolute values are conditional on it: if I instead use Vasavada's 1.48, both
retrieved values move off the searched grid entirely. What survives is the
direction --- Apollo 17 still comes out more conductive.
\end{say}

Be precise, because this is the question most likely to expose over-claiming.
From \texttt{fixed\_input\_sensitivities.json}, at $\chi=1.48$ the Apollo 15
retrieval runs down to the grid floor (1.0) and Apollo 17 runs up to the
ceiling (25.0). Both are \emph{bounds}, not retrievals --- the minimum leaves
the searched range. So the correct statement is: the absolute numbers are
conditional on $\chi=2.7$; the ordering is not.

\begin{trap}
``My result is robust to $\chi$.'' It is not. The \emph{ordering} is robust to
$\chi$; the \emph{magnitudes} are conditional on it. Saying the first thing
plainly is what makes the rest of your uncertainty analysis believable.
\end{trap}

\Q{What are your two boundary conditions?}
\begin{say}
At the top, a non-linear radiative energy balance: absorbed sunlight equals
thermal emission plus whatever conducts into the ground. It has $T_s^4$ in it,
so it is solved by Newton iteration at every time step. At the bottom, a fixed
geothermal flux --- a Neumann condition, not a fixed temperature.
\end{say}
\[
(1-A)\,S=\varepsilon\sigma T_s^{4}+K\,\partial_z T\big|_{z=0},
\qquad
K\,\partial_z T\big|_{z_{\max}}=Q_b
\]
The basal condition is the important one to defend. $Q_b$ is 21 mW m$^{-2}$ at
Apollo 15 and 16 at Apollo 17, from Langseth's revised borehole values.

\Q{Your deep profile keeps rising and never flattens. Isn't that unphysical?}
\begin{say}
It is the physical signature of a real geothermal flux. In periodic steady
state the cycle-mean conductive flux is the same at every depth and equals
$Q_b$, so the mean gradient is $Q_b/K$ --- about 2.3 kelvin per metre. A
profile that flattened would mean zero flux, which contradicts Apollo's own
heat-flow measurement.
\end{say}

Then dispose of the bedrock version of the question. If a high-conductivity
layer sits below the column, it changes $K$ \emph{there}, so it flattens the
gradient \emph{below} it --- but it cannot change the flux passing
\emph{through} the sensor zone. My retrieval reads the gradient at 0.8--2.3 m,
which is set by $Q_b$ and $K$ in that range. Bedrock at 5 m leaves $K_d^{*}$
untouched. The model base sits more than 2.5 m below the deepest datum for
exactly this reason.

\Q{What is the thermal skin depth, and why does it govern your design?}
\begin{say}
It is the depth at which the daily temperature wave has decayed by a factor of
$e$ --- about 5 cm here. It matters because it separates the column into two
regimes: a thin surface skin where the temperature is genuinely
time-dependent, and everything below it where only the cycle mean survives. My
solver exploits exactly that split.
\end{say}

The diurnal amplitude at 0.55 m is roughly 0.15 K --- four orders of magnitude
below the surface swing. That is why the anchor can sit there.

\Q{Why a geometric grid rather than a uniform one?}
\begin{say}
Because the temperature curvature is enormous at the surface and almost zero at
depth. Cells start at 2 mm and grow 8 percent each, reaching 5 m in 69 cells. A
uniform grid fine enough to resolve the skin would need tens of thousands of
cells to reach 5 m, and would spend nearly all of them where nothing is
happening.
\end{say}

Related check question: face conductivities use the \emph{harmonic} mean, not
the arithmetic mean, because two cells in series are two thermal resistances in
series. An arithmetic mean would let a high-conductivity cell wrongly dominate
its low-conductivity neighbour.

% =====================================================================
\section{C \; The anchor method --- your contribution}

\Q{State your methodological contribution in one sentence.}
\begin{say}
Instead of time-stepping the whole five-metre column for thousands of
lunations until it settles, I time-step only the thin sun-baked skin and
reconstruct everything below it by integrating a single ordinary differential
equation downward from an anchor point at 0.55 m.
\end{say}

\pointat{Slides 9 and 10, and the guidebook \texttt{anchormethod\_slide}.}

\Q{What property of the converged state makes it legal to skip the deep column?}
\begin{say}
Average the heat equation over one full cycle. In periodic steady state the
stored-heat term averages to exactly zero --- the cycle closes, so the
temperature ends where it began. Time drops out of the problem, and what
remains says the cycle-mean conductive flux is constant with depth and equal to
$Q_b$. That turns a partial differential equation in depth and time into an
ordinary differential equation in depth alone.
\end{say}

This is the single most important thing to be able to say cleanly. The five
moves are: average over one lunation; the time term dies because the cycle is
periodic; the average commutes with the depth derivative; pin the constant to
$Q_b$ at the base; then peel the average apart, which names the gap
$u_{\rm rect}$.
\[
\frac{d\langle T\rangle}{dz}=\frac{Q_b-u_{\rm rect}}{K(\langle T\rangle,z)}
\]

\begin{trap}
``Because it is steady state.'' That restates the question. You must name the
mechanism: the stored-heat term averages to zero over a periodic cycle, so the
mean flux is depth-independent.
\end{trap}

\pointat{Guidebook \texttt{pde2odeflow} and \texttt{closureflow}.}

\Q{What is $u_{\rm rect}$, and why can you ignore it below the anchor?}
\begin{say}
It is the rectified, or eddy, flux --- the correlation between the fluctuating
part of the conductivity and the fluctuating part of the gradient. Because $K$
depends on temperature, hot moments are also higher-conductivity moments, so
the average of the product is not the product of the averages.
$u_{\rm rect}$ is that difference. It scales as the square of the temperature
fluctuation, so it lives in the skin and dies with it.
\end{say}
\[
u_{\rm rect}\equiv\langle K\,\partial_z T\rangle-K(\langle T\rangle)\,\partial_z\langle T\rangle
\]
At the anchor it is under 1 percent of $Q_b$. Note the sign: the closure is
$(Q_b-u_{\rm rect})/K$, with a \textbf{minus}. Below the anchor the correction
is negligible and the slope is effectively $Q_b/K$.

\Q{Why is the anchor at 0.55 m specifically?}
\begin{say}
Deep enough that the diurnal signal is dead --- about 0.15 K of amplitude, with
$u_{\rm rect}$ below 1 percent of $Q_b$ --- but shallow enough that the skin
solve stays cheap. It is not a guess: the solver runs a two-stage schedule,
settling first at 0.25 m with a loose tolerance and then at 0.55 m with a tight
one, and the choice is certified by a flux-closure check rather than assumed.
\end{say}

\Q{What was the ``F1 bug'', and why does it matter that you found it?}
\begin{say}
The literature convention is to spin up for a fixed number of lunations --- 30
is typical --- and declare convergence when the cycle-to-cycle change drops
below a tolerance. That test passes long before the profile has actually
arrived. The skin settles in about ten lunations while the deep column is still
tens of kelvin away, because the deep column relaxes on a much slower diffusive
clock. Cycle-to-cycle stability is necessary but nowhere near sufficient.
\end{say}

So convergence is tested on \emph{anchor drift} against a tolerance
(0.005 K), plus a flux-closure check, plus an honest 120-lunation re-run from
the converged profile. Three independent checks, because any one of them alone
can be fooled.

\pointat{Guidebook \texttt{certflow}; thesis Fig 3.2.}

\Q{Why 96 inner lunations? The literature would use 12.}
\begin{say}
Because at 12 the answer is still moving. I ran a convergence scan and watched
$K_d^{*}$ itself, not just the flux closure. At Apollo 17, 12 inner lunations
biases the retrieved value high by about 0.37 mW; the value only plateaus in
the 48 to 96 range. Apollo 15 is the slower site --- lower conductivity means
lower diffusivity --- so Apollo 15 is what sets the requirement.
\end{say}

\begin{figure}[htbp]\centering
\includegraphics[width=\linewidth]{rev_convergence.pdf}
{\footnotesize From \texttt{convergence\_scan.json}. Left: the
retrieved value against inner lunations, plotted as the error from the
converged answer, with the $\pm0.005$ mW certification band shaded. Right: the
flux-closure error over the same scan. Both flatten only well past 12.}
\end{figure}

The general lesson, and worth saying out loud: a number that ``looks settled''
earns an explicit re-derivation, not a pass.

\Q{Crank--Nicolson is unconditionally stable. So why is $\Delta t$ 1800 s and not 3600 s?}
\begin{say}
Stability is not the constraint --- accuracy is. Unconditional stability means
the solution will not blow up, not that it is correct. At 3600 s the converged
periodic state was first-order accurate in the time step and carried a
depth-uniform bias, because the last step of each cycle was never taken; the
cycle advanced by one step less than a full period, so the phase slipped.
\end{say}

The fix has two parts: halve the step to 1800 s, and add a \emph{wrap step}
that closes the cycle exactly so the final state is genuinely the next cycle's
starting state. That restored second-order behaviour and moved
$K_d^{*}$(A17) from 7.16 to 7.08. The step is now certified by a $\Delta t$
ladder with a Richardson check, not asserted.

\begin{trap}
``Crank--Nicolson is unconditionally stable so any step is fine.'' This is the
exact reasoning that produced the bug.
\end{trap}

\Q{How much faster, and how do you know it is the same answer?}
\begin{say}
One anchored solve takes about 13 seconds against roughly 260 seconds for a
brute-force solve to comparable accuracy --- about 20 times per solve
algorithmically, and about 2500 times once compiled, so a full sweep goes from
several hours to under a minute. Agreement is not assumed: the anchored and
brute-force solutions agree to better than a hundredth of a milliwatt.
\end{say}

Always pair the speed claim with the accuracy claim in the same breath. A
speed-up with no equivalence check is worthless to a committee, and inviting
the question is worse than answering it unprompted.

\Q{Why not just use a faster computer, or write it in C++?}
\begin{say}
Both, and neither is the point. The solver is already compiled --- it runs
through Numba, and there is an independent C++ implementation that reproduces
it to about twelve significant figures as a cross-check. But hardware buys you
a constant factor. The anchor method removes the work rather than doing it
faster, which is why it scales to a global model of millions of columns where
buying 2500 times more computer does not.
\end{say}

% =====================================================================
\section{D \; Retrieval, statistics and model selection}

\Q{How is $K_d$ actually found?}
\begin{say}
By a direct sweep, not an optimiser. I try roughly thirty candidate values, run
a full certified solve for each, and compute the root-mean-square misfit
against the deep sensors. That traces a bowl; I fit a parabola through the
three points bracketing the minimum and take the vertex as the retrieved value.
\end{say}

The grid is deliberately non-uniform, with a fine band bracketing the expected
minimum, because a coarse uniform sweep plus a three-point parabola carries
vertex noise of a couple of hundredths of a milliwatt --- which was once enough
to make two runs disagree in the second decimal.

\Q{Why a brute-force sweep rather than gradient descent or least squares?}
\begin{say}
Three reasons. The forward model is expensive and non-linear, and each
evaluation is a full converged solve, so gradients are not free. A sweep gives
me the entire misfit curve, not just the minimum, which is what lets me see
whether the minimum is sharp or flat. And once the anchor method made each
solve cost a second, the sweep became cheap enough that there was no reason to
be clever.
\end{say}

\Q{What is the bootstrap doing, and why is it the right tool here?}
\begin{say}
It answers: how much of my answer is an accident of which sensors I happened to
have? Each of 1500 draws resamples the deep sensors with replacement and
jitters every depth by the real emplacement uncertainty, plus or minus about
2.5 cm, then re-fits. The spread of the 1500 answers is the statistical
uncertainty the data can support.
\end{say}

Two details worth having: the draws are \emph{paired} across sites, so the
contrast is computed draw by draw rather than by differencing two independent
distributions; and the re-fit reads cached sweep profiles rather than
re-solving, which is what makes 1500 draws affordable. The seed is archived.

\Q{You quote 0.031. Say exactly what that number is.}
\begin{say}
It is a bootstrap tail proportion: of the 1500 paired draws, 46 produced a
contrast at or below zero --- that is, Apollo 17 not coming out more
conductive. So $P_{\rm boot}(\Delta K_d\le 0)=0.031$. It is not a $p$-value.
\end{say}

A $p$-value is computed under an assumed null hypothesis --- a world in which
the two sites are genuinely identical. I never constructed that world. What I
did was resample my own sensors and re-run my own retrieval. Calling the result
a $p$-value would import a hypothesis test I did not perform.

\begin{trap}
``$p<0.05$, so the difference is significant.'' Two errors in five words: it is
not a $p$-value, and the two-sided interval contains zero. Say ``bootstrap tail
proportion'' every time. Your guidebook makes the same correction.
\end{trap}

\Q{Then how can 0.031 be less than 0.05 while your 95\,\% interval contains zero?}
\begin{say}
Because they are two different cuts of the same distribution. The tail
proportion is one-sided. The confidence interval is two-sided --- it removes
2.5 percent from each end. For zero to fall outside a two-sided 95 percent
interval, the lower tail would have to be smaller than 0.025. Mine is 0.031.
Since 0.031 is greater than 0.025, zero is inside the interval necessarily.
There is no tension between the two numbers at all.
\end{say}

\begin{figure}[htbp]\centering
\includegraphics[width=0.92\linewidth]{rev_tails.pdf}
{\footnotesize The same 1500 paired draws, cut two ways. Contrast
median 2.31, 95\,\% interval $[-0.12, 3.56]$.}
\end{figure}

If you can draw this on a whiteboard in ten seconds, this question stops being
dangerous and becomes the moment you look most in command of your own
statistics.

\Q{What is AICc, and what does it tell you here?}
\begin{say}
It asks whether an extra fitted parameter pays for itself. A better fit always
lowers the residual sum of squares, so AICc adds a penalty for each parameter,
with a small-sample correction that bites hard when you have few data points.
I compare three models at each site: the global value with nothing fitted, the
Hayne form with $K_d$ fitted per site, and a three-layer piecewise model.
\end{say}

The result, from \texttt{model\_selection.json}, with $\Delta$ computed as
fitted minus global:

\vspace{4pt}
{\small
\begin{tabular}{@{}lccccc@{}}
\toprule
site & $N$ & global & fitted & 3-layer & $\Delta$(fit $-$ global) \\
\midrule
Apollo 15 & 7  & 4.06  & 7.00   & 6.71 & \textbf{$+2.94$} \\
Apollo 17 & 16 & $-1.39$ & $-24.56$ & 2.40 & \textbf{$-23.17$} \\
\bottomrule
\end{tabular}}
\vspace{4pt}

Decisive at Apollo 17. \textbf{Not} decisive at Apollo 15 --- there the penalty
exceeds the fit improvement.

\Q{So Apollo 15's own model selection does not justify fitting it separately. Why do you still quote 4.60?}
\begin{say}
Because AICc asks whether the extra parameter pays for itself \emph{in that
dataset}, and with 7 sensors spanning only 84 to 139 cm, it does not. That is a
statement about the power of the Apollo 15 data, not evidence that the site
actually has the global value. Failing to reject is not accepting. And the 95
percent interval at Apollo 15 is 4.18 to 6.96, which excludes 3.4.
\end{say}

Then close it deliberately: the two tools answer different questions, I report
both, and the thesis claim rests on the interval and on the ordering rather
than on Apollo 15's AICc. Volunteering this before you are pushed is worth more
than defending it after.

\pointat{Backup: thesis Fig 5.6, the AICc decomposition.}

\Q{Why is a three-layer model in the comparison, and what did it show?}
\begin{say}
As a control. If a discrete layered structure fitted the deep sensors better
than the smooth Hayne form, that would undermine using Hayne's functional form
at all. It does not win at either site. At Apollo 17 its retrieved value is
also pinned at the grid floor, so that row is a bound rather than a fair
comparison, and I say so.
\end{say}

\Q{What does the MCMC add that the bootstrap does not?}
\begin{say}
They answer different questions. The bootstrap holds the basal flux fixed and
asks how much the \emph{data} could move my answer. The MCMC lets $Q_b$ float
under a prior and asks how much the \emph{flux assumption} could move it.
Neither is redundant; one measures noise and the other measures a systematic.
\end{say}

Implementation, if asked: 32 walkers, 4000 steps, first 1000 discarded, giving
96{,}000 retained samples with acceptance around 0.69 and autocorrelation
times of 32--36. The likelihood reads a bilinear interpolation of a directly
solved RMSE surface --- 117 real solves per site on a 13 by 9 grid --- and
rejects anything outside it rather than extrapolating.

\begin{trap}
Do not oversell it. There is no R-hat, no trace plot and no posterior
predictive check in the thesis, and the noise level in the likelihood is
uncalibrated. If asked about convergence diagnostics, say what exists ---
autocorrelation and effective sample size --- and what does not.
\end{trap}

\Q{Your MCMC says the $K_d$--$Q_b$ correlation is $-0.75$ at Apollo 17 but $+0.08$ at Apollo 15. Why so different, and does it help or hurt you?}
\begin{say}
At Apollo 17, sixteen sensors spanning 1.3 to 2.3 m resolve the gradient well
enough to see the trade-off clearly: raise the assumed flux and the retrieved
conductivity must fall to keep the same slope. That is the degeneracy, measured.
At Apollo 15, seven sensors over a 55 cm span cannot tilt the posterior, so it
stays nearly round. It hurts the magnitude and helps the ordering.
\end{say}

The reason it helps the ordering: because the two sites respond to a flux
revision in \emph{opposite} directions, a systematic error in $Q_b$ does not
move them together. Across the whole published flux envelope the probability
that Apollo 17 exceeds Apollo 15 stays at about 99.2 percent, and it
\emph{rises} as the prior widens.

\pointat{Backup: thesis Fig 5.11 and Fig 5.13.}

\Q{Where does your uncertainty actually come from?}
\begin{say}
Not mainly from the statistics. The bootstrap contributes about 0.52 mW at
both sites. The systematic terms are larger, and the dominant one is different
at each site: the basal flux dominates at Apollo 15, and the surface albedo
dominates at Apollo 17. The quadrature totals are 1.88 and 3.88 mW.
\end{say}

\begin{figure}[htbp]\centering
\includegraphics[width=\linewidth]{rev_budget.pdf}
{\footnotesize From \texttt{kd\_error\_budget.json}. Statistical
precision is good; accuracy is limited by the fixed inputs.}
\end{figure}

Say the last line out loud: \emph{statistical precision is good, and accuracy
is limited by the systematics.} That is the honest summary and it pre-empts the
question.

% =====================================================================
\section{E \; Robustness and the honest limits}

\Q{What is your single largest caveat?}
\begin{say}
That my absolute values are conditional on Hayne's radiative coefficient
$\chi=2.7$. Under Vasavada's 1.48 both retrievals leave the searched grid.
What survives that change is the direction of the contrast, not its size. I
state this in the thesis rather than rounding it away.
\end{say}

Deliver this without hedging. A committee is far more worried about a candidate
who has not found their own weakness than about the weakness.

\Q{Is there any prediction your model makes that was never fitted?}
\begin{say}
Yes --- the surface temperature. The retrieval only ever sees deep-sensor
temperatures below 80 cm, so the modelled diurnal surface swing is a genuine
out-of-sample prediction. I compare it against the Diviner orbital composite,
and I should be honest about the outcome: it reproduces the shape and the
amplitude, but the Hayne per-site fit carries a warm bias of about 2.6 K at
Apollo 15 and 4.1 K at Apollo 17.
\end{say}

Two things save this and you should say both. First, the retrieval is driven by
the deep \emph{gradient}, not the surface value, so a surface offset does not
propagate into $K_d$ in the way it would for an orbital fit. Second, the
Mart\'inez--Siegler forward model closes the surface better --- biases of
$+0.08$ and $-0.78$ K --- which is evidence that the residual is in the surface
parameterisation rather than in my deep retrieval.

\begin{trap}
``The Diviner comparison validates the model.'' It is a genuine out-of-sample
test and it partly passes. Overstating it is the fastest way to lose a
committee that has read Fig 5.10.
\end{trap}

\Q{What happens if you hold data back?}
\begin{say}
The probes carry two independent thermometer bridges, so I can fit on one and
predict the other. At Apollo 15 the TG-only fit predicts the TR sensors with
RMSE 1.074 K against an in-sample 1.052 K --- essentially no degradation. I
also drop the deepest sensor and re-fit, which is the harshest single-point
test because the deepest sensor carries the most leverage on the gradient.
\end{say}

\pointat{Backup: thesis Fig 5.7.}

\Q{What would falsify your result?}
\begin{say}
Several things, and I can name them. A revision of the Apollo basal heat flows
that moved both sites the same way and far enough. A demonstration that the
borestem disturbance reaches deeper than 80 cm. A radiative coefficient
substantially different from 2.7 that also changed the sign of the difference
between the two sites --- though every value I tested preserved the ordering.
Or a third borehole that disagreed.
\end{say}

Having a falsification answer ready is a strong signal. A candidate who cannot
name one is usually asked why not.

\Q{You found three bugs in your own work. Why should I trust the fourth one isn't there?}
\begin{say}
You should not trust that --- you should look at what the process does when
there is one. Each of the three was found by a specific check, not by luck: a
wrong basal flux, an under-resolved spin-up, and a time step that was
first-order because the cycle never closed. Each is now covered by a
convergence scan, a ladder, or a certification test that runs every build, and
every input carries a register entry saying whether it is literature-founded,
convergence-founded or budgeted.
\end{say}

The three, with what each cost: $Q_b$(A17) 15 to 16 mW m$^{-2}$ moved
$K_d^{*}$(A17) from 8.12 to 7.51; the spin-up fix moved it 7.51 to 7.16; the
time-step fix moved it 7.16 to 7.08. Note the trend --- each correction made
the answer smaller and the analysis more conservative.

\pointat{Backup: guidebook \texttt{threebugs} and \texttt{inputflow}.}

% =====================================================================
\section{F \; The doctoral plan}

{\color{dim}\small The doctoral Q\&A is 7 minutes --- longer than the thesis
Q\&A. Rehearse this section at least as hard.}

\Q{What is the gap between what you did and what you propose?}
\begin{say}
The thesis measured two points. The question people actually want answered ---
where can subsurface ice survive --- is about the whole Moon. The thesis built
the engine that makes that tractable: a solve now costs about a second, and a
global map is millions of independent columns. The doctorate is about turning
two calibration points into a global product.
\end{say}

\Q{Why start with terrain shading rather than going straight to a global map?}
\begin{say}
Because my thesis assumed flat ground, and neither site is flat. Apollo 15 sits
at the foot of the Apennine front and Apollo 17 in a valley between massifs.
Using real elevation models, the horizon reaches 14.0 degrees at Apollo 15,
removing 1.16 percent of the incoming sunlight, and 10.1 degrees and 0.18
percent at Apollo 17. This is already implemented and was presented at AOGS.
\end{say}

Stress that it is \emph{built}, not proposed. The rhetorical structure of Part
2 is three slides of completed work before any promise.

\Q{Which ground property actually controls the answer?}
\begin{say}
Density, but through conductivity rather than through heat storage. Density does
two jobs --- it sets how much heat the regolith stores and how well it conducts.
I ran 650 model variations to separate them. Letting density set the
conductivity as well as the storage cuts the misfit at Apollo 15 from 1.74 to
0.90 kelvin, roughly halving it. At Apollo 17 it moves very little, 0.39 to
0.36, because that site already fits well.
\end{say}

\begin{trap}
``Heat capacity barely matters and conductivity matters a lot.'' That is true
at Apollo 15 and not at Apollo 17. State both. A committee member who checks
the AOGS numbers will find the asymmetry, and finding it after you have
overstated is much worse than you having said it.
\end{trap}

\Q{You claim the layered physics is transferable. But moving it across doubles the error.}
\begin{say}
The comparison that carries the argument is not crossed against native --- it
is crossed against uniform. A uniform-property ground gives 2.31 kelvin at
Apollo 15 and 3.76 at Apollo 17. Taking Apollo 15's best configuration,
untouched, and applying it at Apollo 17 gives 1.73 --- still more than twice as
good as the uniform assumption, from a configuration tuned at a completely
different site. If this were curve-fitting to one hole, moving it would buy
nothing.
\end{say}

Then concede the rest on purpose: yes, it degrades from 0.36 to 1.73, and that
degradation is exactly why the doctorate proposes to \emph{learn} the property
field rather than assume one configuration is universal. The gap is the
motivation, not a weakness.

\Q{Is a Moon-wide model computationally feasible?}
\begin{say}
Yes, and that is the point of the thesis method. The columns are independent,
so the problem is embarrassingly parallel. At about a second per converged
column, a global map at a practical resolution is a large but ordinary
computation --- and it was simply not reachable at 27 hours per column, which
is why nobody had done it.
\end{say}

Have a number ready if pressed, and be honest that the exact resolution is a
trade-off to be set in year two rather than something already fixed.

\Q{What are the three deliverables, one per year?}
\begin{say}
Year one, a validated property model: finish the density-linked conductivity
work, go beyond the standard functional form, and test against both boreholes.
Year two, a Moon-wide subsurface temperature map with real terrain shading,
checked against orbital measurements. Year three, an ice-survivability map
built from those profiles, delivered to the sounding missions that need it.
\end{say}

Say each deliverable out loud. That is the line a committee listens for.

\Q{What is the largest risk in year one, and what is your mitigation?}
\begin{say}
That the property model does not generalise beyond the two boreholes --- that
each site simply wants its own parameters. My mitigation is the cross-site test
I already ran: I check transferability before extrapolating, not after. If the
cross-applied configuration stopped beating a uniform ground, that would be the
signal to change approach rather than push on.
\end{say}

\Q{What is the largest risk in year two?}
\begin{say}
Validation coverage. A global map is only as good as what I can check it
against, and the only in-situ anchors are the two boreholes. The mitigation is
to validate against orbital surface temperatures everywhere and against the
boreholes at depth, and to be explicit that the deep field is a physical
extrapolation constrained at two points rather than a measurement.
\end{say}

Saying that plainly is better than being asked. The honest framing is that the
product is a physically-constrained model with two ground-truth anchors, which
is more than any current global thermal model has.

\Q{What if your global model disagrees with Diviner?}
\begin{say}
That would be informative rather than fatal, and I already have a version of
this result. My thesis retrieval carries a warm surface bias against Diviner
under the Hayne form while the Mart\'inez--Siegler form closes better. A
disagreement at the surface points at the surface parameterisation; a
disagreement at depth would point at the property model. The two are separable
because they are constrained by different data.
\end{say}

\Q{Why does this need a doctorate rather than being more of the same?}
\begin{say}
The thesis answered a closed question at two points with a fixed functional
form. The doctorate has to do three things it did not: replace an assumed
property model with a learned one, carry real three-dimensional terrain into a
one-dimensional column model in a physically defensible way, and turn a
temperature field into an ice-stability product that a mission can actually
use. Each of those is a research problem, not an extension.
\end{say}

\Q{Why here, and why now?}
\begin{say}
The solver exists and is verified, the terrain and property studies are
finished and were presented at AOGS, and the laboratory has the link to the
NICT terahertz sounding mission that needs exactly this boundary condition.
The subsurface-sounding instruments are being built now; the temperature
profiles they need do not exist yet.
\end{say}

\Q{Suppose you finish and the answer is that the Moon really is uniform after all. Was it worth it?}
\begin{say}
Yes, and I would publish that. It would be the first time the uniformity
assumption had been tested rather than inherited, and every ice-stability
estimate that rests on it would gain a justification it currently does not
have. My thesis already suggests otherwise at both boreholes, but the value of
the work is in making the assumption testable, not in which way it falls.
\end{say}

\vspace{18pt}
\hrule
\vspace{8pt}
{\footnotesize\color{dim}
Built from \texttt{code/results/*.json} on the certified run. Figures:
\texttt{rev\_tails}, \texttt{rev\_convergence} and \texttt{rev\_budget} are
generated by \texttt{make\_reviewer\_art.py}; all other figures referenced are
the thesis and guidebook originals.}

\end{document}
